
\begin{definition}
    Пусть $M$ -- гладкое ориентированное многообразие, $K \subset M$ -- компакт.
    Пусть $\omega \in \Omega^m(M)$ и $\operatorname{supp} \omega \subset K$.
    Определим интеграл по компакту $K$:
    \[ \int_K \omega := \int_M \rho \omega, \]
    где $\rho \in C_c^\infty(M)$ -- такая функция, что $\rho(x) = 1$ для всех $x \in K$.
    Существование $\rho$ следует из теоремы о разбиении единицы.
    Корректность следует из того, что если $\rho_1|_K = \rho_2|_K = 1$, то носитель $(\rho_1 - \rho_2)$ лежит вне $K$, где $\omega \equiv 0$. Следовательно, $(\rho_1 - \rho_2)\omega \equiv 0$ на $M$.
\end{definition}

\subsection{Теорема Стокса}

Операторы внешнего дифференцирования $d$ форм и взятия края $\partial$ связаны.

\begin{theorem}[\textbf{Стокс}]
    Пусть $M$ -- гладкое $m$-мерное ориентируемое многообразие ($m > 1$), $N$ -- регулярная область в $M$. Тогда для любой $\omega \in \Omega_c^{m-1}(M)$ выполнено
    \[
    \int_N d\omega = \int_{\partial N} i^*\omega,
    \]
    где ориентации $N$ и $\partial N$ согласованы.
\end{theorem}

\begin{myproof}
	Сначала установим справедливость формулы для полупространства, а затем распространим результат на произвольное многообразие с помощью разбиения единицы.
\begin{lemma}
    Теорема верна для $M = \R^m$ и $N = \mathbb{H}^m$.
\end{lemma}

\begin{prooflemma}
    Пусть форма $\omega$ имеет вид $\omega = \sum_{i=1}^m f_i(x)\, dx_1 \wedge \dotsc \wedge \widehat{dx_i} \wedge \dotsc \wedge dx_m$. Тогда:
    %($\widehat{dx_i}$ означает, что $dx_i$ пропускается при записи). 
    \[
    d\omega = \sum_{i,j=1}^m \frac{\partial f_i}{\partial x_j}(x)\, dx_j \wedge dx_1 \wedge \dotsc \wedge \widehat{dx_i} \wedge \dotsc \wedge dx_m = \sum_{i=1}^m (-1)^{i-1} \frac{\partial f_i}{\partial x_i}(x)\, dx_1 \wedge \dotsc \wedge dx_m.
    \]
    Поэтому $\int_{\mathbb{H}^m} d\omega = \sum_{i=1}^m (-1)^{i-1} \int_{\mathbb{H}^m} \frac{\partial f_i}{\partial x_i}(x)\, dx_1 \dots dx_m$.
    При $i > 1$ по теореме Фубини имеем:
    \[
    \int_{\mathbb{H}^m} \frac{\partial f_i}{\partial x_i}(x)\, dx_1 \dots dx_m = \int_{(-\infty, 0) \times \R^{m-2}} \left( \int_{-\infty}^\infty \frac{\partial f_i}{\partial x_i}(x)\, dx_i \right) dx_1 \dots \widehat{dx_i} \dots dx_m.
    \]
    По формуле Ньютона-Лейбница, учитывая компактность $\supp(f_i)$, имеем:
    \[
    \int_{-\infty}^\infty \frac{\partial f_i}{\partial x_i}(x)\, dx_i = f_i(x_1, \dots, x_i, \dots, x_m) \Big|_{x_i=-\infty}^{x_i=\infty} = 0.
    \]
    Следовательно, от суммы останется только слагаемое при $i=1$, то есть:
    \[
    \int_{\mathbb{H}^m} d\omega = \int_{\R^{m-1}} \left( \int_{-\infty}^0 \frac{\partial f_1}{\partial x_1}(x)\, dx_1 \right) dx_2 \dots dx_m = \int_{\R^{m-1}} f_1(0, x_2, \dots, x_m)\, dx_2 \dots dx_m.
    \]
	С другой стороны, вложение $i_0 : \R^{m-1} \cong \partial \mathbb{H}^m \to \R^m$ имеет вид $i_0(x') = (0, x')$,
	поэтому $i_0^*\omega = f_1(0, x_2, \dots, x_m)\, dx_2 \wedge \dotsc \wedge dx_m$, остальные слагаемые занулятся. Значит,
    \[
    \int_{\R^{m-1}} i_0^*\omega = \int_{\R^{m-1}} f_1(0, x_2, \dots, x_m)\, dx_2 \dots dx_m.
    \]
    Полученный интеграл совпадает с $\int_{\partial \mathbb{H}^m} i_0^* \omega$, так как ориентация $\partial \mathbb{H}^m$ для стандартной ориентации $\mathbb{H}^m$ соответствует $(x_2, \dots, x_m)$. Требуемое равенство установлено.
\end{prooflemma}

\begin{note}
    Если вместо $\mathbb{H}^m$ в доказательстве леммы рассматривать $\R^m$, то не надо исключать случай $i=1$, так что $\int_{\R^m} d\omega = 0$.
\end{note}

Докажем теорему.
Для каждой точки $p \in \supp(\omega)$ найдется параметризация $\phi: V \to U_p$ с условием: если $p \in \supp(\omega) \cap \partial N$, то $\phi(V \cap \mathbb{H}^m) = U_p \cap N$, а если $p \in \supp(\omega) \setminus \partial N$, то $U_p$ не пересекается с $\partial N$.
Уменьшая окрестности, можно считать их связными. Также считаем, что все полученные параметризации согласованы с ориентацией, иначе заменим на композицию с линейным отображением.
Из открытого покрытия $\{U_p\}$ компакта $\supp(\omega)$ выделим конечное подпокрытие.
Получим набор параметризаций $\phi_i: V_i \to U_i$ ($1 \leq i \leq n$), покрывающих $\supp(\omega)$ и согласованных с ориентацией.
Пусть $\{\rho_i\}_{i=1}^n$ -- разбиение единицы, подчиненное этому покрытию. Поскольку $\omega = \sum_{i=1}^n \rho_i \omega$ и $d\omega = d(\sum_{i=1}^n \rho_i \omega) = \sum_{i=1}^n d(\rho_i \omega)$,
то в силу линейности интеграла достаточно доказать формулу в предположении, что носитель $\omega$ покрывается образом одной параметризации $\phi: V \to U$. Возможны случаи.

\vspace{5pt}
Случай $U \cap \partial N \ne \emptyset$. Рассмотрим сужение $\phi_0: V_0 \to U_0$ параметризации $\phi$ в терминологии леммы 8.1, а также вложения $i: \partial N \to M$ и $i_0: \R^{m-1} \cong \partial \mathbb{H}^m \to \R^m$. 
Тогда $i \circ \phi_0 = \phi \circ i_0$. Распишем левую и правую часть формулы.
\[
\int_N d\omega = \int_{V \cap \mathbb{H}^m} \phi^*(d\omega) = \int_{\mathbb{H}^m} \phi^*(d\omega) = \int_{\mathbb{H}^m} d(\phi^*\omega),
\]
\[
\int_{\partial N} i^*\omega = \int_{V_0} \phi_0^*(i^*\omega) = \int_{\R^{m-1}} \phi_0^*(i^*\omega) = \int_{\R^{m-1}} i_0^*(\phi^*\omega) = \int_{\partial \mathbb{H}^m} i_0^*(\phi^*\omega).
\]
По лемме последние интегралы в строках совпадают, что доказывает формулу.

\vspace{5pt}
Случай $U \subset \text{int}\, N$. Имеем $\int_N d\omega = \int_V \phi^*(d\omega) = \int_{\R^m} \phi^*(d\omega) = \int_{\R^m} d(\phi^*\omega) = 0$ по замечанию после леммы. 
Поскольку перенос $i^*\omega = 0$, то формула верна и в этом случае.

\vspace{5pt}
Случай $U \subset \text{ext}\, N$. Сужение $d \omega$ на $N$ и перенос $i^* \omega$ нулевые, поэтому формула также верна. Все возможные случаи разобраны, доказательство завершено.
\end{myproof}

\begin{corollary}
    Пусть $M$ -- гладкое $m$-мерное ориентируемое многообразие, $\omega \in \Omega_c^{m-1}(M)$. 
	Тогда $\int_M d\omega = 0$, так как $\partial M = \emptyset$.
\end{corollary}

Опишем согласованность ориентаций области и ее края еще одним способом.

\begin{definition}
    Пусть $N$ -- регулярная область в $M$, $p \in \partial N$ и $v \in T_p M$. 
    Пусть $\varphi: V \to U$ -- параметризация окрестности $p$ ($\varphi(a)=p$), выпрямляющая границу, то есть $\varphi(V \cap \mathbb{H}^m) = U \cap N$. 
    Рассмотрим вектор $w = (d\varphi_a)^{-1}(v) = (w_1, \dots, w_m) \in \R^m$. 
    Тогда:
    \begin{enumerate}[itemsep = 0pt, topsep = 5pt]
        \item Если $w_1 = 0$, то $v \in T_p(\partial N)$ (касательный вектор),
        \item Если $w_1 > 0$, то $v$ называется \textit{внешним} вектором,
        \item Если $w_1 < 0$, то $v$ называется \textit{внутренним} вектором.
    \end{enumerate}
\end{definition}

\begin{note}
    Определение корректно, так как любая допустимая замена координат $\Phi$ отображает полупространство $x_1 < 0$ в себя. Следовательно, ось $x_1$ не меняет направления и знак первой координаты вектора $w$ сохраняется.
\end{note}

\begin{proposition}
    Пусть $N$ -- регулярная область в ориентируемом многообразии $M$. Ориентации $N$ и $\partial N$ согласованы $\Lra$ для всякой точки $p \in \partial N$ из условия, что базис $(v_1, \dots, v_{m-1})$ положительно ориентирован в $T_p(\partial N)$ следует, что базис $(n, v_1, \dots, v_{m-1})$ положительно ориентирован в $T_p M$.
\end{proposition}

Запишем частные случаи формулы Стокса в $\R^m$ для $m=2$ и $m=3$.

\begin{example}[\textbf{формула Грина}]
Пусть $G$ -- ограниченная регулярная область в $\R^2$, причем ориентация края $\partial G$ согласована со стандартной (индуцированной с $\R^2 = \langle e_1, e_2 \rangle$) ориентацией $G$. Пусть $\omega = P\, dx + Q\, dy \in \Omega^1(\R^2)$. Тогда
\[
\int_{\partial G} P\, dx + Q\, dy = \iint_G \left( \frac{\partial Q}{\partial x} - \frac{\partial P}{\partial y} \right) dx \, dy.
\]
Согласованность ориентаций $G$ и $\partial G$ в данном случае означает, что $\partial G$ ориентирована внешней нормалью $n$.
Неформально говоря, интегрирование ведется в таком направлении обхода контура, при котором область $G$ остается слева. 
\end{example}

\begin{example}[\textbf{формула Гаусса-Остроградского}]
Пусть $G$ -- ограниченная регулярная область в $\R^3$, причем ориентация $\partial G$ согласована со стандартной (индуцированной с $\R^3$) ориентацией $G$. Пусть $\omega = P\, dy \wedge dz + Q\, dz \wedge dx + R\, dx \wedge dy \in \Omega^2(\R^3)$. Тогда
\[
\int_{\partial G} P\, dy \wedge dz + Q\, dz \wedge dx + R\, dx \wedge dy = \iiint_G \left( \frac{\partial P}{\partial x} + \frac{\partial Q}{\partial y} + \frac{\partial R}{\partial z} \right) dx \, dy \, dz.
\]
Согласованность ориентаций $G$ и $\partial G$ в данном случае означает то же, что и выше.
\end{example}

\begin{example}[\textbf{классическая формула Стокса}]
Пусть $M$ -- гладкое 2-мерное ориентируемое многообразие в $\R^3$, а $S$ -- ограниченная регулярная область в $M$, покрываемая одной параметризацией $r: V \to \R^3$. Пусть $\omega = P\, dx + Q\, dy + R\, dz \in \Omega^1(\R^3)$. Тогда
\[
\int_{\partial S} P\, dx + Q\, dy + R\, dz = \int_S \left( \frac{\partial R}{\partial y} - \frac{\partial Q}{\partial z} \right) dy \wedge dz + \left( \frac{\partial P}{\partial z} - \frac{\partial R}{\partial x} \right) dz \wedge dx + \left( \frac{\partial Q}{\partial x} - \frac{\partial P}{\partial y} \right) dx \wedge dy.
\]
Согласованность ориентаций $S$ и $\partial S$ в данном случае означает следующее: пусть $p \in \partial S$, $N(p)$ -- единичная нормаль к $S$ (задающая ориентацию $S$), $\tau(p)$ -- касательный вектор к $\partial S$ (задающий ориентацию края), а $n(p)$ -- вектор внешней нормали к краю $\partial S$, лежащий в касательной плоскости $T_p M$. Тогда базис $(n, \tau, N)$ положительно ориентирован в $\R^3$.
\end{example}

\subsection{Замкнутые и точные дифференциальные формы}

\begin{definition}
	Пусть $M$ -- гладкое многообразие, $\omega \in \Omega^{k\geq 1}(M)$.
    Форма $\omega$ \textit{замкнута}, если $d\omega = 0$.
    Форма $\omega$ \textit{точна}, если $\omega = d\alpha$ для некоторой формы $\alpha \in \Omega^{k-1}(M)$.
\end{definition}

\begin{note}
	Из условия $d^2 = 0$ следует, что всякая точная форма замкнута.
\end{note}

\begin{theorem}[\textbf{лемма Пуанкаре}]
    Если форма $\omega \in \Omega^{k\geq 1}(\R^n)$ замкнута, то она точна.
\end{theorem}

\begin{myproof}
    Пусть $(t, x) = (t, x_2, \dots, x_n)$ -- координаты в $\R^n$.
	Отметим, что всякая $k$-форма на $\R^n$ есть сумма мономов вида $f(t, x) \cdot dt \wedge dx^I$, где $I \in \mathbb{I}_{k-1}$, и $g(t, x)\cdot dx^J$, где $J \in \mathbb{I}_k$.
    
	\vspace{5pt}
    Рассмотрим линейный оператор $\Phi$, который на мономах действует следующим образом: если $\omega = f(t, x) \cdot dt \wedge dx^I$, то $\Phi(\omega) = \left( \int_0^t f(s, x) ds \right) dx^I$, и если $\omega = g(t, x) \cdot dx^J$, то $\Phi(\omega) = 0$. 
	По сути оператор $\Phi$ действует как <<взятие первообразной по координате $t$>>.
    
	\vspace{5pt}
    Покажем, что значение $\Phi$ на всякой $\omega \in \Omega^k(\R^n)$ удовлетворяет соотношению:
    \[
    d\Phi(\omega) + \Phi(d\omega) = \omega - \pi^*(i^*\omega), \eqno(*)
    \]
    где $\pi: \R^n \to \R^{n-1}$ -- проектирование ($\pi(t, x) = x$), а $i: \R^{n-1} \to \R^n$ -- вложение ($i(x) = (0, x)$).
    В силу линейности соотношение ($*$) достаточно проверить для мономов.
    
	\vspace{5pt}
    Пусть $\omega = g(t, x) \cdot dx^J$. Первое слагаемое в левой части ($*$) равно нулю по определению $\Phi$. Заметим, что $d\omega = \frac{\partial g}{\partial t} \cdot dt \wedge dx^J + \omega_0$, где форма $\omega_0$ не содержит $dt$, поэтому
    \[
    \Phi(d\omega) = \left( \int_0^t \frac{\partial g}{\partial s}(s, x) ds \right) dx^J = g(t, x) \cdot dx^J - g(0, x)\cdot dx^J.
    \]
    Поскольку $(i \circ \pi)(t, x) = (0, x)$, то $g(0, x) \cdot dx^J = (i \circ \pi)^*(\omega)$ и равенство ($*$) установлено.
    
	\vspace{5pt}
    Пусть $\omega = f(t, x) \cdot dt \wedge dx^I$. Тогда $\Phi(\omega) = \left( \int_0^t f(s, x) ds \right) dx^I$, а значит,
    \[
    d\Phi(\omega) = f(t, x) \cdot dt \wedge dx^I + \sum_{i=2}^n \left( \int_0^t \frac{\partial f}{\partial x_i}(s, x) ds \right) dx_i \wedge dx^I =: \omega + \omega_1.
    \]
    Далее, $d\omega = \sum_{i=2}^n \frac{\partial f}{\partial x_i}(t, x) dx_i \wedge dt \wedge dx^I$, поэтому $\Phi(d\omega) = -\omega_1$. Итак, $d\Phi(\omega) + \Phi(d\omega) = \omega$ и $(i \circ \pi)^*(\omega) = 0$, поэтому равенство ($*$) установлено и в этом случае.
    
	\vspace{5pt}
    Полученное равенство ($*$) позволяет доказать утверждение индукцией по $n$. Для базы индукции предположим, что $n=k$. Если $k$-форма $\omega$ замкнута в $\R^k$,
	то $d\omega = 0$ и $i^*(\omega) = 0$, поэтому ($*$) примет вид $d\Phi(\omega) = \omega$. Это доказывает точность формы $\omega$. 
	Для шага индукции предположим, что всякая замкнутая $k$-форма в $\R^n$ при $n \ge k$ точна. Пусть $k$-форма $\omega$ замкнута в $\R^{n+1}$. 
	Тогда $di^*(\omega) = i^*(d\omega) = 0$, то есть $i^*(\omega)$ замкнута в $\R^n$. По предположению существует форма $\alpha$ в $\R^n$, что $i^*(\omega) = d\alpha$. Поэтому из ($*$) имеем $\omega = d\Phi(\omega) + \pi^*(d\alpha) = d\Phi(\omega) + d(\pi^*\alpha) = d(\Phi(\omega) + \pi^*\alpha)$, так что форма $\omega$ точна.
\end{myproof}